%%%%%%%%%%%%%%%%%%%%%%%%%%%%%%%%%%%%%%%%%%%%%%%%%%%%%%%%%%%%%%%%%%%%%%
% Lớp 12 | Các số đặc trưng ... ghép nhóm | Phương sai, độ lệch chuẩn của mẫu số liệu ghép nhóm
% Dạng 2D3?2-3: Câu hỏi tổng hợp
% Số câu: 12  (sinh tự động bởi tach_cau.py; từ nay thầy sửa trực tiếp file này)
%%%%%%%%%%%%%%%%%%%%%%%%%%%%%%%%%%%%%%%%%%%%%%%%%%%%%%%%%%%%%%%%%%%%%%

%%%%%%%%%%%%%%%% TRẮC NGHIỆM 4 PHƯƠNG ÁN %%%%%%%%%%%%%%%%

% nguồn: Lop_12_Dai_so_Chuong_3
\begin{ex}%[2D3H2-3]%[ma:ae524ae3]
	Bảng dưới đây thống kê độ ẩm không khí trung bình mỗi ngày đo được tại một trạm  quan trắc trong 60 ngày đầu năm 2023.
	\begin{center}
		\resizebox{0.8\linewidth}{!}{\begin{tikzpicture}
		\draw[opacity=.25,thin,step=.2,cyan](0,0) grid  (7,2.5);
		 \draw[opacity=1,cyan](0,0) grid (7,2.5);
		\draw[stealth-stealth](0,2.5) node[left]{Tần số tương đối (\%)}|-(7,0)  node[below,anchor=west]{Độ ẩm (\%)};
		\foreach\x/\dientich[count=\i from  1] in {0/60,1/64,1/68,1.5/72,1/76,0.5/80}{
			\draw[pattern=north west lines,pattern  color=cyan!65](\i-1,0) rectangle +(1,\x);
			 \draw (\i,0) node[below]{\dientich};
		 }
		 \foreach \y [count=\i from 1] in {10,20,30}{
			\draw (-.1,\i/2)--(.1,\i/2) 
			(0,\i/2)  node[left]{$\y$};}
		\foreach \z [count=\i from 1]  in {20,20,30,20,10}{
			\draw  (\i+0.5,\z/20) node[above] {$\z$};}
	\end{tikzpicture}
		 }
	\end{center}
	Số trung bình của mẫu số liệu ghép  nhóm trên là
	\choice 
		 {$68{,}2$}
		{$69$} 
		 {\True $69{,}2$}
		 {$70$}
	\loigiai{Ta có bảng thống  kê theo giá trị đại diện như sau
		\begin{center}
			\begin{tabular}{|c|c|c|c|c|c|c|}
				\hline 
				Gía trị đại diện&  $62$ & $66$ & $70$ & $74$ & $78$ \\
				 \hline
				Tần số &$12$& $12$ &  $18$ &$12$ & $6$ \\
				 \hline
			\end{tabular}
		\end{center}
		Cỡ mẫu là $n=60$. Số trung bình của mẫu số  liệu ghép nhóm trên là \[\overline{x}=\dfrac{12\cdot 62+12\cdot 66+18\cdot 70+12\cdot 74+6\cdot 78}{60}=69{,}2.\]
	 }
\end{ex}

% nguồn: Lop_12_Dai_so_Chuong_3
\begin{ex}%[2D3H2-3]%[ma:f117e108]
	Điểm trung bình môn  Toán của một số học sinh khối $12$ được ghi lại ở bảng sau
	\begin{center}
		\begin{tabular}{|l|c|c|c|c|c|}
			\hline Nhóm  &{$[5;6)$}&{$[6;7)$}&{$[7;8)$}&{$[8;9)$}&{$[9;10)$}\\
			\hline Tần số & $10$  & $14$ & $19$ & $32$ & $25$ \\
			 \hline
		\end{tabular}
	\end{center}
	Số trung bình của mẫu  số liệu ghép nhóm trên là
	 \choice
		 {$7{,}5$}
		{$7$} 
		 {\True $7{,}98$}
		 {$8{,}65$}
	\loigiai{Ta có bảng thống kê  theo giá trị đại diện như sau
		\begin{center}
			\begin{tabular}{|c|c|c|c|c|c|c|}
				 \hline
				Điểm& $5{,}5$ & $6{,}5$ & $7{,}5$ & $8{,}5$ & $9{,}5$ \\ 
				 \hline
				Tần số &$10$& $14$ & $19$ &$32$ & $25$ \\ 
				 \hline
			\end{tabular}
		\end{center}
		Tổng số học sinh là $n=100$. Số  trung bình của mẫu số liệu ghép nhóm trên là \[\overline{x}=\dfrac{10\cdot 5{,}5+14\cdot 6{,}5+19\cdot 7{,}5+32\cdot 8{,}5+25\cdot 9{,}5}{100}=7{,}87.\]
	} 
\end{ex}


%%%%%%%%%%%%%%%% ĐÚNG/SAI %%%%%%%%%%%%%%%%

% nguồn: Lop_12_Dai_so_Chuong_3
\begin{ex}%[2D3H2-3]%[ma:1f3ffa0d]
	Kết quả $ 40$  lần nhảy xa của hai vận động viên Dũng và Huy được lần lượt thống kê trong bảng bên dưới (đơn vị: mét)
	\begin{longtable}{|c|c|c|}
		\hline 
		Nhóm & Dũng  & Huy\\
		 \hline
		$[6{,}22;6{,}46)$ &  $3$ & $2$\\
		\hline 
		$[6{,}46;6{,}70)$ & $7$  & $5$\\
		 \hline
		$[6{,}70;6{,}94)$  & $5$ & $8$\\
		\hline 
		$[6{,}94;7{,}18)$ & $20$ &  $19$\\
		\hline 
		$[7{,}18;7{,}42)$ & $5$ & $6$\\ 
		\hline 
		& $n=40$ &  $n=40$\\
		\hline 
	\end{longtable}
	Xét tính đúng  sai của các khẳng định sau
	 \choiceTF
		{\True Số trung bình cộng của mẫu số  liệu ghép nhóm biểu diễn kết quả $40$ lần nhảy xa của vận động viên Dũng cho bởi Bảng 15 (làm tròn kết quả đến hàng phần trăm) là $ 6{,}92$ m}
		{\True Độ lệch chuẩn của mẫu số liệu ghép  nhóm biểu diễn kết quả $40$ lần nhảy xa của vận động viên Dũng cho bởi Bảng 15 (làm tròn kết quả đến hàng phần trăm) là $ 0{,}26$ m}
		{Phương sai của mẫu số liệu ghép nhóm biểu diễn kết quả $40$ lần nhảy xa của vận động viên Huy cho bởi Bảng 16 (làm  tròn kết quả đến hàng phần trăm) là $ 0{,}16$}
		{Kết quả nhảy xa của vận động viên Dũng  đồng đều hơn kết quả nhảy xa của vận động viên Huy}
	\loigiai{ 
		\begin{itemchoice}
			\itemch Số trung bình  cộng của mẫu số liệu ghép nhóm biểu diễn kết quả 40 lần nhảy xa của vận động viên Dũng là\\
			$\overline{x}_D=\dfrac{3\cdot 6{,}34+7\cdot 6{,}58+5\cdot 6{,}82+20\cdot  7{,}06+5\cdot 7{,}30}{40}=\dfrac{276{,}88}{40}\approx 6{,}92$(m).
			\itemch Phương sai của mẫu số liệu ghép nhóm biểu diễn kết quả  40 lần nhảy xa của vận động viên Dũng (làm tròn kết quả đến hàng phần trăm) là\\
			$\begin{aligned}
			& s_D^2=\dfrac{1}{40}\left[3\cdot \left(6{,}34-6{,}92\right)^2+7\cdot  \left(6{,}58-6{,}92\right)^2+5\cdot \left(6{,}82-6{,}92\right)^2\right.\\
			&\left.+20\cdot \left(7{,}06-6{,}92\right)^2+5\cdot \left(7{,}30-6{,}92\right)^2\right]=\dfrac{2{,}9824}{40}\approx  0{,}07.\\
		\end{aligned}$\\
		Độ lệch chuẩn của mẫu số liệu ghép  nhóm trên là $s_D\approx\sqrt{0{,}07}\approx 0{,}26$ (m).
		\itemch Số trung bình cộng của mẫu số liệu ghép nhóm biểu diễn kết quả  40 lần nhảy xa của vận động viên Huy là\\
		$\overline{x}_H=\dfrac{2\cdot 6{,}34+5\cdot 6{,}58+8\cdot 6{,}82+19\cdot 7{,}06+6\cdot 7{,}30}{40}=\dfrac{278{,}08}{40}\approx 6{,}95$(m)\\ 
		Vậy phương sai của mẫu số liệu ghép nhóm biểu diễn kết quả 40 lần nhảy xa của vận động viên Huy (làm tròn kết quả đến hàng  phần trăm) là\\
		$ s_H^2=\dfrac{1}{40}\left[2\cdot \left(6{,}34-6{,}95\right)^2+5\cdot \left(6{,}58-6{,}95\right)^2+8\cdot  \left(6{,}82-6{,}95\right)^2+19\cdot \left(7{,}06-6{,}95\right)^2\right.$\\
		 $\left.+6\cdot \left(7{,}30-6{,}95\right)^2\right]=\dfrac{2{,}5288}{40}\approx 0{,}06$.\\
		Độ lệch chuẩn  của mẫu số liệu ghép nhóm trên là $s_H\approx\sqrt{0{,}06}\approx 0{,}24$ (m).
		\itemch Do $s_H\approx 0{,}24<s_D\approx 0{,}26$ nên kết quả nhảy xa của vận động viên Huy đồng đều hơn kết quả nhảy xa của vận động viên Dũng. 
	\end{itemchoice}
} 
\end{ex}

% nguồn: Lop_12_Dai_so_Chuong_3
\begin{ex}%[2D3H2-3]%[ma:f0fd3cd9]
	Thống kê nhiệt độ trung bình của tháng $2$ tại thành phố Thanh Hóa từ $2\,004$ đến  hết $2\,023$ ($20$ năm) được kết quả sau
	\begin{center}
		\begin{tabular}{|c|c|}
			\hline 
			Lớp nhiệt độ ($^\circ C$)& Tần số \\ 
			 \hline
			$\left[12;14\right)$ & $1$ \\ 
			 \hline
			$\left[14;16\right)$  & $3$ \\
			 \hline
			$\left[16;18\right)$ &  $12$ \\
			 \hline
			 $\left[18;20\right)$ & $9$ \\
			\hline 
			$\left[20;22\right)$ &  $5$ \\
			\hline 
			 Tổng & $30$ \\
			 \hline
		\end{tabular}
	\end{center}
	Khi đó 
	\choiceTF 
		{\True Khoảng biến thiên của bảng  số liệu là~$10$ ($^\circ C$)}
		{Giá trị trung bình là~$18$($^\circ C$)  (làm tròn đến hàng đơn vị)}
		{Phương sai của mẫu số liệu này  là~$3{,}94$ ($^\circ C$) (làm tròn đến hàng phần trăm)}
		{\True Độ lệch chuẩn của mẫu số liệu này là~$1{,}98$  ($^\circ C$) (làm tròn đến hàng phần trăm)}
	 \loigiai{
		\begin{itemchoice}
			\itemch Khoảng biến  thiên của mẫu số liệu
			\centerline {$R=$ Giá trị lớn nhất-Giá trị nhỏ nhất  $=22-12=10(^\circ C)$.}
			\itemch Ta có  số trung bình của mẫu số liệu
			$$\overline{x}=\dfrac{1\cdot  13+15\cdot 3+12\cdot 17+9\cdot 19+5\cdot 21}{30} \approx 17{,}93(^\circ C).$$
			\itemch Phương sai của mẫu số liệu  là
			\[s_x^2=\dfrac{1\cdot (13-17{,}93)^2+3\cdot (15-17{,}93)^2+12\cdot (17-17{,}93)^2+9\cdot  (19-17{,}93)^2+5\cdot (21-17{,}93)^2}{30} \approx 3{,}93(^\circ C) \].
			\itemch Độ lệch chuẩn của mẫu số liệu $s_x=\sqrt{s_x^2} \approx 1{,}98(^\circ C)  $.
		\end{itemchoice}
	} 
\end{ex}

% nguồn: Lop_12_Dai_so_Chuong_3
\begin{ex}%[2D3H2-3]%[ma:e8e1d2d9]
	Mẫu số liệu ghép nhóm thống kê mức  lương của hai công ty $A$, $B$ (đơn vị: triệu đồng) được thể hiện như bảng dưới đây
	\begin{longtable}{|c|c|c|}
		\hline 
		Nhóm & Công ty A & Công ty B\\ 
		\hline 
		$[10;15)$ & $15$  & $25$\\
		 \hline
		 $[15;20)$ & $18$ & $15$\\
		\hline 
		 $[20;25)$ & $10$ & $7$\\
		\hline 
		$[25;30)$ & $10$ & $5$\\ 
		 \hline
		$[30;35)$  & $5$ & $5$\\
		\hline 
		$[35;40)$  & $2$ & $3$\\
		\hline 
		 & $n=60$ & $n=60$\\
		\hline 
	\end{longtable}
	Xét tính đúng sai  của các khẳng định sau
	 \choiceTF
		{\True Số trung bình cộng của mẫu số liệu ghép nhóm của công  ty A là $\dfrac{62}{3}$}
		{Độ lệch chuẩn của mẫu số liệu ghép nhóm của công ty A là  $ 5$}
		{Phương sai của mẫu số liệu ghép  nhóm của công ty B (làm tròn kết quả đến hàng phần trăm) là $ 52{,}91$}
		{\True Nhận thấy độ lệch chuẩn của công ty A nhỏ hơn công ty B nên mức lương của công ty  A đồng đều hơn}
	 \loigiai{
		\begin{itemchoice}
			\itemch Số trung bình cộng của mẫu  số liệu ghép nhóm của công ty A là\\
			$\overline{x}_A=\dfrac{15\cdot 12{,}5+18\cdot 17{,}5+10\cdot 22{,}5+10\cdot 27{,}5+5\cdot 32{,}5+2\cdot  37{,}5}{60}=\dfrac{62}{3}$.
			\itemch Phương sai của mẫu số liệu ghép nhóm của công  ty $A$ là\\
			$\begin{aligned}
			 &{s_A}^2=\dfrac{15\cdot{\left(12{,}5-\dfrac{62}{3}\right)^2}+18\cdot{\left(17{,}5-\dfrac{62}{3}\right)^2}+10\cdot{\left(22{,}5-\dfrac{62}{3}\right)^2}+10\cdot{\left(27{,}5-\dfrac{62}{3}\right)^2}}{60}\\
			&+\dfrac{5\cdot{\left(32{,}5-\dfrac{62}{3}\right)^2}+2\cdot{\left(37{,}5-\dfrac{62}{3}\right)^2}}{60}\approx 49{,}14.\\ 
		\end{aligned}$\\
		Độ lệch chuẩn của mẫu số liệu ghép nhóm của công ty  $A$ là $s_A=\sqrt{s_A^2}=\sqrt{49{,}13}\approx 7$.
		\itemch Số trung bình  cộng của mẫu số liệu ghép nhóm của công ty $B$ là\\
		$\overline{x}_B=\dfrac{25\cdot  12{,}5+15\cdot 17{,}5+7\cdot 22{,}5+5\cdot 27{,}5+5\cdot 32{,}5+3\cdot 37{,}5}{60}=\dfrac{229}{12}$\\
		Phương sai của mẫu số  liệu ghép nhóm của công ty $B$ là\\
		$\begin{aligned}
		&{s_B}^2=\dfrac{25\cdot \left(12{,}5-\dfrac{229}{12}\right)^2+15\cdot \left(17{,}5-\dfrac{229}{12}\right)^2+7\cdot \left(22{,}5-\dfrac{229}{12}\right)^2+5\cdot  \left(27{,}5-\dfrac{229}{12}\right)^2}{60}\\
		&+\dfrac{5\cdot \left(32{,}5-\dfrac{229}{12}\right)^2+3\cdot \left(37{,}5-\dfrac{229}{12}\right)^2}{60}\approx  57{,}91.\\
	\end{aligned}$
	\itemch  Độ lệch chuẩn của mẫu số liệu ghép nhóm của công ty B là $s_B=\sqrt{s_B^2}=\sqrt{57{,}91}\approx 7{,}61$.\\
	Nhận  thấy độ lệch chuẩn của công ty $A$ nhỏ hơn công ty $B$ nên mức lương của công ty $A$ đồng đều hơn.
	\end{itemchoice}
 }
\end{ex}

% nguồn: Lop_12_Dai_so_Chuong_3
\begin{ex}%[2D3H2-3]%[ma:f86068d2]
	Biểu đồ sau mô tả kết quả điều tra về điểm  trung bình năm học của học sinh hai trường $A$ và $B$
	\begin{center}
		\begin{tikzpicture}[scale=1,font=\footnotesize,line join=round,line cap=round,>=stealth]
				\begin{scope}[scale=0.8]
					\draw[->] (0,0)--(12,0) node [below  right]{Điểm trung bình};
					\draw[->] (0,0)--(0,8)  node[above left]{Số học sinh};
					\foreach \y in  {1,2,3,4,5,6,7} \draw (1pt,\y)--(-1pt,\y) node[left]{$\y$};
					\foreach  \x/\y in{0/4,2/5,4/3,6/4,8/2} \draw[pattern=dots] (\x,\y) rectangle (\x+1,0);
					\foreach \x/\y in{1/2,3/5,5/4,7/3,9/1} \draw[pattern=north west lines]  (\x,\y) rectangle (\x+1,0);
					\node  at (1,0) [below]{$[5; 6)$};
					\node at (3,0) [below]{$[6; 7)$}; 
					\node  at (5,0) [below]{$[7; 8)$};
					\node at (7,0)  [below]{$[8; 9)$};
					\node at  (9,0) [below]{$[9; 10)$};
					\draw[pattern=dots] (9,7) rectangle (9.3,7.3) node[above  right]{Học sinh trường A};
					\draw[pattern=north west lines]  (9,6) rectangle (9.3,6.3) node[above right]{Học sinh trường B};
				\end{scope}
			\end{tikzpicture}
	\end{center}
	Người ta lập được bảng tần số ghép nhóm cho  mẫu số liệu trên như sau
	\begin{longtable}{|c|c|c|c|c|c|}
		\hline 
		Điểm trung bình & $[5;6)$ & $[6;7)$ & $[7;8)$ & $[8;9)$ & $[9;10)$\\ 
		 \hline
		Giá trị đại diện & $5{,}5$ & $6{,}5$ & $7{,}5$ &  $8{,}5$ & $9{,}5$\\
		 \hline
		Học sinh trường A & $4$  & $5$ & $3$ & $4$ & $2$\\
		 \hline
		Học sinh trường B & $2$ & $5$ & $4$  & $3$ & $1$\\
		 \hline
	\end{longtable}
	Xét tính đúng sai của các khẳng  định sau
	 \choiceTF
		{\True Tứ phân vị  thứ nhất của mẫu số liệu ghép nhóm của học sinh trường A là $ 6{,}1$}
		{\True Khoảng tứ phân vị của mẫu  số liệu ghép nhóm của học sinh trường B là $ 1{,}73$}
		{\True Nếu so sánh theo khoảng tứ phân vị của mẫu số liệu ghép nhóm thì học sinh trường B có  điểm trung bình đồng đều hơn}
		{Nếu so sánh theo độ lệch chuẩn của mẫu số liệu ghép nhóm thì học sinh trường A có điểm trung bình đồng đều  hơn}
	 \loigiai{
		\begin{itemchoice}
			\itemch Cỡ mẫu $n_A=18$\\ 
			Gọi $x_1$; $x_2$; $\ldots$; $x_{18}$ là  mẫu số liệu gốc về điểm trung bình năm học của học sinh hai trường $A$ được xếp theo thứ tự không giảm.\\
			Ta có: $x_1$; $\ldots$; $x_4\in\left[5;6\right)$; $x_5$;$\ldots$; $x_9\in\left[6;7\right)$; $x_{10}$;$\ldots$;$x_{12}\in\left[7;8\right)$; $x_{12}$;$\ldots$;$x_{16}\in\left[8;9\right)$; $x_{17}$; $x_{18}\in\left[9;10\right)$.\\ 
			Tứ phân vị thứ nhất của mẫu số liệu gốc là  $x_5\in\left[6;7\right)$.\\
			Do đó, tứ phân vị thứ nhất của mẫu  số liệu ghép nhóm là $Q_1=6+\dfrac{\dfrac{18}{4}-4}{5}\left(7-6\right)=6{,}1$\\
			Tứ phân vị thứ ba của  mẫu số liệu gốc là $x_{14}\in\left[8;9\right)$.\\
			Do đó, tứ phân  vị thứ ba của mẫu số liệu ghép nhóm là $$Q_3=8+\dfrac{\dfrac{3\cdot 18}{4}-\left(4+5+3\right)}{4}\left(9-8\right)=8{,}375.$$
			Khoảng tứ phân vị  của mẫu số liệu ghép nhóm là $\Delta_Q=Q_3-Q_1=2{,}275$
			 \itemch Cỡ mẫu $n_B=15$\\
			Gọi $y_1$; $y_2$; $\ldots$; $y_{15}$ là mẫu số liệu gốc về điểm trung  bình năm học của học sinh hai trường $B$ được xếp theo thứ tự không giảm.\\
			Ta có: $y_1$; $y_2\in\left[5;6\right)$;  $y_3$;$\ldots$; $y_7\in\left[6;7\right)$; $y_8$; $\ldots$; $y_{11}\in\left[7;8\right)$; $y_{12}$; $\ldots$; $y_{14}\in\left[8;9\right)$; $y_{15}\in\left[9;10\right)$.\\
			Tứ phân vị thứ nhất của mẫu số liệu gốc là $y_4\in\left[6;7\right)$.\\ 
			Do đó, tứ phân vị thứ nhất  của mẫu số liệu ghép nhóm là $Q'_1=6+\dfrac{\dfrac{15}{4}-2}{5}\left(7-6\right)=6{,}35$.\\
			Tứ phân vị thứ ba của mẫu số liệu  gốc là $y_{12}\in\left[8;9\right)$.\\
			Do đó  tứ phân vị thứ ba của mẫu số liệu ghép nhóm là $$Q'_3=8+\dfrac{\dfrac{3\cdot 15}{4}-\left(2+5+4\right)}{3}\left(9-8\right)=8{,}08.$$
			Khoảng tứ phân vị  của mẫu số liệu ghép nhóm là $\Delta'_Q=Q'_3-Q'_1=1{,}73$.
			\itemch Vậy nếu so sánh theo khoảng  tứ phân vị của mẫu số liệu ghép nhóm thì học sinh trường B có điểm trung bình đồng đều hơn.
			\itemch Xét số  liệu của trường $A$:\\
			 Số trung bình $\overline{x}_A=\dfrac{4\cdot 5{,}5+5\cdot 6{,}5+3\cdot 7{,}5+4\cdot 8{,}5+2\cdot 9{,}5}{18}=7{,}22$.\\
			Độ lệch chuẩn $s_A=\sqrt{\dfrac{4\cdot 5{,}5^2+5\cdot 6{,}5^2+3\cdot 7{,}5^2+4\cdot  8{,}5^2+2\cdot 9{,}5^2}{18}-7{,}22^2}\approx 1{,}79$.\\
			 Xét số liệu của trường $B$:\\
			Số trung  bình $\overline{x}_B=\dfrac{2\cdot 5{,}5+5\cdot 6{,}5+4\cdot 7{,}5+3\cdot 8{,}5+1\cdot 9{,}5}{15}=7{,}23$.\\
			Độ  lệch chuẩn $s_B=\sqrt{\dfrac{2\cdot 5{,}5^2+5\cdot 6{,}5^2+4\cdot 7{,}5^2+3\cdot 8{,}5^2+1\cdot 9{,}5^2}{15}-7{,}23^2}\approx 1{,}31$.\\
			Vậy nếu so sánh theo độ lệch chuẩn của mẫu số liệu ghép nhóm thì học sinh trường B có điểm trung bình  đồng đều hơn.
		\end{itemchoice}
	} 
\end{ex}

% nguồn: Lop_12_Dai_so_Chuong_3
\begin{ex}%[2D3V2-3]%[ma:35c9a7a1]
	Cô Phương thống kê lại số giờ chơi thể thao trong 1 tuần của  học sinh lớp $10C$ ở bảng sau
	\begin{center}
		\begin{tabular}{|l|c|c|c|c|}
			\hline 
			\textbf{Số giờ} & $\left[0;3 \right)$&$\left[3;6 \right)$&$\left[6;9  \right)$&$\left[9;12 \right)$ \\
			 \hline
			\textbf{Số học sinh} & $3$ &  $10$ & $14$ &$23$\\
			\hline 
		\end{tabular}
	\end{center}
	 \choiceTF
		 {\True Khoảng biến thiên của mẫu số liệu ghép nhóm trên là 12 (giờ)}
		{\True  Tứ phân vị thứ nhất của mẫu số liệu gốc thuộc $\left[3;6 \right)$}
		{Khoảng tứ  phân vị của mẫu số liệu ghép nhóm trên là $\dfrac{681}{460}$}
		{\True Phương sai của mẫu số liệu ghép  nhóm trên là $7{,}9236$}
	\loigiai{ 
		\begin{itemchoice}
			\itemch {\bf Đúng.}\\ 
			Khoảng biến thiên của mẫu số liệu ghép nhóm trên  là $12-0=12$ (giờ).
			 \itemch {\bf Đúng.}\\
			Cỡ mẫu $n=50$.\\ 
			Gọi $x_1;x_2;x_3;...;x_{50}$ là mẫu số liệu gốc gồm thời gian chơi thể thao  trong 1 tuần của học sinh lớp $10C$ được xếp theo thứ tự không giảm.\\
			Ta có $x_1;x_2;x_3\in \left[0;3 \right)$; $x_4;...;x_{13}\in \left[3;6  \right)$; $x_{14};...;x_{27}\in \left[6;9 \right)$; $x_{28};...;x_{50}\in \left[9;12 \right)$.\\
			Tứ phân vị thứ  nhất của mẫu số liệu gốc là $x_{13}\in \left[3;6 \right)$.
			\itemch {\bf Sai.}\\ 
			Tứ phân vị thứ nhất của mẫu số liệu ghép nhóm là \[Q_1=3+\dfrac{\dfrac{50}{4}-3}{10}\cdot \left( 6-3  \right)=\dfrac{117}{20}.\]
			Tứ phân vị thứ ba  của mẫu số liệu gốc là $x_{38}\in \left[9;12 \right)$.\\
			Do đó tứ phân vị thứ  ba của mẫu số liệu ghép nhóm là \[Q_3=9+\dfrac{\dfrac{3\cdot 50}{4}-\left( 3+10+14 \right)}{23}\cdot \left( 12-9 \right)=\dfrac{477}{46}.\]
			Vậy khoảng tứ phân vị của mẫu số liệu ghép nhóm  là \[\Delta_Q=Q_3-Q_1=\dfrac{477}{46}-\dfrac{117}{20}=\dfrac{2079}{460}.\]
			\itemch {\bf  Đúng.}\\
			Ta có  bảng thống kê số giờ chơi thể thao trong 1 tuần của học sinh lớp $10C$ theo giá trị đại diện:
			\begin{center}
				\begin{tabular}{|l|c|c|c|c|}
					 \hline
					\textbf{Số giờ}  & $\left[0;3 \right)$&$\left[3;6 \right)$&$\left[6;9 \right)$&$\left[9;12 \right)$ \\
					 \hline
					\textbf{Giá trị đại diện} & $1{,}5$ &  $4{,}5$ & $7{,}5$ & $10{,}5$ \\
					 \hline
					\textbf{Số học sinh} & $3$ & $10$ & $14$ &$23$\\ 
					\hline 
				\end{tabular}
			\end{center}
			Số trung bình của mẫu số liệu ghép nhóm là \[\overline{x}=\dfrac{3\cdot  1{,}5+10\cdot 4{,}5+14\cdot 7{,}5+23\cdot 10{,}5}{50}=7{,}92.\]
			 Phương sai của mẫu số liệu ghép nhóm trên là
			\[S^2=\dfrac{3\cdot  1{,}5^2+10\cdot 4{,}5^2+14\cdot 7{,}5^2+23\cdot 10{,}5^2}{50}-7{,}92^2=7{,}9236.\]
		\end{itemchoice}
	 }
\end{ex}


%%%%%%%%%%%%%%%% TRẢ LỜI NGẮN %%%%%%%%%%%%%%%%

% nguồn: Lop_12_Dai_so_Chuong_3
\begin{bt}%[2D3N2-3]%[ma:fe56635a]
	Bảng dưới đây biểu diễn mẫu số liệu ghép nhóm về cản nặng của một số quả bưởi da xanh ở một khu vực (đơn vị: kg). 
	\begin{center}
		\begin{tabular}{|c|c|c|c|c|c|}
			 \hline
			Nhóm&$\left[1{,}2;1{,}3\right)$ &$\left[1{,}3;1{,}4\right)$ &$\left[1{,}4;1{,}5\right)$ &  $\left[1{,}5;1{,}6\right)$ & $\left[1{,}6;1{,}7\right)$ \\
			\hline 
			Tần số&8 &21 &  8 & 7 &6 \\
			 \hline
		\end{tabular}
	\end{center}
	Số trung bình của mẫu số liệu ghép nhóm trên là bao nhiêu? (Làm  tròn kết quả đến hàng phần trăm).
	\shortans{1,41} 
	\loigiai{ 
		Số trung bình của mẫu số liệu ghép nhóm là 
		\[\overline{x}=\dfrac{8\cdot 1{,}25+21\cdot 1{,}35+8\cdot 1{,}45+7\cdot 1{,}55+6\cdot  1{,}65}{50}\approx 1{,}41.\]
	} 
\end{bt}

% nguồn: Lop_12_Dai_so_Chuong_3
\begin{bt}%[2D3H2-3]%[ma:7e49d6d5]
	Thống kê điểm thi cuối kì  $1$ môn Toán của lớp $10$A$1$ ta được bảng sau
	\begin{center}
		\begin{tabular}{|c|c|c|c|c|c|c|c|c|c|}
			\hline 
			Điểm & $2$ &  $5$ & $5{,}5$ & $6$ & $7$ & $8$ & $8{,}5$ & $9$ & $10$ \\ \hline
			Số học sinh & $1$ & $2$ & $3$ & $9$ & $11$ & $13$ &  $5$ & $2$ & $1$ \\ \hline
		\end{tabular}
	\end{center}
	Hãy cho biết mẫu số  liệu trên có bao nhiêu giá trị ngoại lệ?
	 \shortans{1}
	 \loigiai{
		Ta có  $Q_2=M_e=x_{24}=7$; $Q_1=x_{12}=6$; $Q_3=x_{36}=8$ và $\Delta_Q=Q_3-Q_1=2$.\\
		Gọi $x$ là giá trị  ngoại lệ, khi đó $\hoac{&x>Q_3+1{,}5\Delta_Q \\&x<Q_1-1{,}5\Delta_Q}
		\Leftrightarrow \hoac{&x>8+1{,}5\cdot 2  \\&x<6-1,5\cdot 2}
		 \Leftrightarrow \hoac{&x>11 \\&x<3.}$\\
		Vậy  có đúng một giá trị ngoại lệ là $2$.
	 }
\end{bt}


%%%%%%%%%%%%%%%% TỰ LUẬN %%%%%%%%%%%%%%%%

% nguồn: Lop_12_Dai_so_Chuong_3
\begin{bt}%[2D3H2-3]%[ma:f444f1f2]
	Biểu đồ dưới đây mô tả kết quả điều tra về mức lương khởi điểm (đơn vị: triệu đồng)  của một số công nhân ở hai khu vực $A$ và $B$.
	\begin{enumerate}
		\item Hãy xác định  giá trị đại diện cho mỗi nhóm và lập bảng tần số ghép nhóm cho mẫu số liệu đó.
		\item Nếu so sánh theo độ  lệch chuẩn của mẫu số liệu ghép nhóm thì công nhân ở khu vực nào có mức lương khởi điểm đồng đều hơn?
	\end{enumerate}
	Mức  lương khởi điểm của công nhân ở hai khu vực $A$ và $B$
	\begin{center}
		\begin{tikzpicture}[scale=1,>=stealth,font=\footnotesize,line join=round,line cap=round,blue,x=7mm,y=7mm]
				\def\xmin{0} \def\xmax{13} 
				\def\ymin{0}  \def\ymax{9}
				\foreach \d in {1,2,...,8}{\draw[green,ultra thin]  (0,\d)--++(\xmax,0);}
				\draw[->] (\xmin,0)--(\xmax,0) node [below]{Mức  lương};
				\draw[->]  (0,\ymin)--(0,\ymax) node [left]{Số công nhân};
				\foreach \y in  {0,1,...,8}
				\draw (1pt,\y)--(-1.5pt,\y) node [left] {$\y$}; 
				\fill[color=yellow,opacity=.75] (1,0)--(1,4)--(2,4)--(2,0); 
				 \fill[color=orange,opacity=.75] (2,0)--(2,3)--(3,3)--(3,0);
				 \fill[color=yellow,opacity=.75] (3,0)--(3,5)--(4,5)--(4,0);
				\fill[color=orange,opacity=.75]  (4,0)--(4,6)--(5,6)--(5,0);
				 \fill[color=yellow,opacity=.75] (5,0)--(5,5)--(6,5)--(6,0);
				\fill[color=orange,opacity=.75] (6,0)--(6,5)--(7,5)--(7,0); 
				 \fill[color=yellow,opacity=.75] (7,0)--(7,4)--(8,4)--(8,0);
				 \fill[color=orange,opacity=.75] (8,0)--(8,5)--(9,5)--(9,0);
				 \fill[color=yellow,opacity=.75] (9,0)--(9,2)--(10,2)--(10,0);
				 \fill[color=orange,opacity=.75] (10,0)--(10,1)--(11,1)--(11,0);
				 \fill[color=yellow,opacity=.75] (9,7)--(9,7.5)--(9.5,7.5)--(9.5,7) node[right,yshift=1mm]{Khu vực A};
				\fill[color=orange,opacity=.75]  (9,6)--(9,6.5)--(9.5,6.5)--(9.5,6) node[right,yshift=1mm]{Khu vực B};
				\path (-1,9.5) node[right]{\bf Mức lương khởi điểm của công  nhân ở hai khu vực A và B};
				\path (2,0) node[below]{$[5;  6)$}
					(4,0) node[below]{$[6;  7)$}
					(6,0) node[below]{$[7; 8)$} 
					(8,0) node[below]{$[8;  9)$}
					(10,0) node[below]{$[9; 10)$} 
					(\xmax,0)  node[below=3mm]{(triệu đồng)};
			\end{tikzpicture}
	\end{center}
	\loigiai{ 
		\begin{enumerate}
			 \item Ta có bảng sau
			\begin{longtable}{|c|c|c|c|c|c|}
				\hline 
				Mức lương &  $[5;6)$ & $[6;7)$ & $[7;8)$ & $[8;9)$ & $[9;10)$\\
				\hline 
				Mức lương đại diện (triệu đồng) & $5{,}5$  & $6{,}5$ & $7{,}5$ & $8{,}5$ & $9{,}5$\\
				 \hline
				 Khu vực A & $4$ & $5$ & $5$ & $4$ & $2$\\
				\hline 
				Khu vực B & $3$ & $6$ & $5$ &  $5$ & $1$\\
				\hline 
			\end{longtable}
			\item Xét mẫu số liệu  của khu vực $A$ :\\
			Cỡ  mẫu là $n_A=4+5+5+4+2=20$.\\
			Số trung bình của mẫu số liệu ghép nhóm là  $$\overline{x}_A=\dfrac{4\cdot 5{,}5+5\cdot 6{,}5+5\cdot 7{,}5+4\cdot 8{,}5+2\cdot 9{,}5}{20}=7{,}25.$$
			 Phương sai của mẫu số liệu ghép nhóm là\\
			$$  S_A^2=\dfrac{1}{20}\left(4\cdot{5{,}5^2}+5\cdot{6{,}5^2}+5\cdot{7{,}5^2}+4\cdot{8{,}5^2}+2\cdot{9{,}5^2}\right)-(7{,}25)^2=1{,}5875.$$
			Độ lệch chuẩn của mẫu số liệu ghép nhóm là  $S_A=\sqrt{1{,}5875}$.\\
			 Xét mẫu số liệu của khu vực $B$ :\\
			Cỡ mẫu là $n_B=3+6+5+5+1=20$.\\ 
			Số trung bình của mẫu số liệu ghép nhóm là $$\overline{x}_B=\dfrac{3\cdot 5{,}5+6\cdot  6{,}5+5\cdot 7{,}5+5\cdot 8{,}5+1\cdot 9{,}5}{20}=7{,}25.$$
			Phương  sai của mẫu số liệu ghép nhóm là\\
			 $$ S_B^2=\dfrac{1}{20}\left(3\cdot{5{,}5^2}+6\cdot{6{,}5^2}+5\cdot{7{,}5^2}+5\cdot{8{,}5^2}+1\cdot{9{,}5^2}\right)-(7{,}25)^2=1{,}2875.$$
			Độ lệch chuẩn của mẫu số liệu ghép nhóm là $S_B=\sqrt{1{,}2875}$.\\ 
			Do $S_A>S_B$ nên nếu so sánh theo độ lệch chuẩn của mẫu số liệu ghép nhóm thì mức lương khởi điểm của công nhân  khu vực $B$ đồng đều hơn của công nhân khu vực $A$.
		\end{enumerate}
	 }
\end{bt}

% nguồn: Lop_12_Dai_so_Chuong_3
\begin{bt}%[2D3H2-3]%[ma:ab42a6b6]
	Giá đóng  cửa của một cổ phiếu là giá của cổ phiếu đó cuối một phiên giao dịch. Bảng sau thống kê giá đóng cửa (đơn vị: nghìn đồng) của hai mã cổ phiếu $A$ và $B$ trong 50 ngày giao dịch liên tiếp.
	\begin{longtable}{|c|c|c|c|c|c|}
		 \hline
		Giá đóng cửa &  $[120;122)$ & $[122;124)$ & $[124;126)$ & $[126;128)$ & $[128;130)$\\
		\hline 
		Cổ phiếu A & $8$ & $9$ & $12$ & $10$ & $11$\\ 
		 \hline
		Cổ phiếu B & $16$ &  $4$ & $3$ & $6$ & $21$\\
		 \hline
	\end{longtable}
	Người ta có thể dùng phương sai và độ lệch chuẩn để so sánh mức độ rủi ro của các loại cổ phiếu có giá trị trung bình gần bằng nhau. Cổ phiếu nào có phương sai, độ lệch chuẩn cao hơn thì  được coi là có độ rủi ro lớn hơn.\\
	Theo quan  điểm trên, hãy so sánh độ rủi ro của cổ phiếu $A$ và cổ phiếu $B$.
	 \loigiai{
		Ta có bảng thống kê  giá đóng cửa theo giá trị đại diện
		\begin{longtable}{|c|c|c|c|c|c|}
			 \hline
			Giá đóng cửa &  $121$ & $123$ & $125$ & $127$ & $129$\\
			\hline 
			Cổ  phiếu A & $8$ & $9$ & $12$ & $10$ & $11$\\
			\hline 
			Cổ phiếu  B & $16$ & $4$ & $3$ & $6$ & $21$\\
			\hline 
		\end{longtable}
		 Xét mẫu số liệu của cổ phiếu $A$\\
		Số trung bình của mẫu số liệu ghép  nhóm là
		$$\overline{x}_1=\dfrac{8\cdot 121+9\cdot 123+12\cdot 125+10\cdot 127+11\cdot 129}{50}=125{,}28.$$ 
		Phương  sai của mẫu số liệu ghép nhóm là
		$$S_1^2=\dfrac{1}{50}\left(8\cdot 121^2+9\cdot 123^2+12\cdot 125^2+10\cdot  127^2+11\cdot 129^2\right)-(125{,}28)^2=7{,}5216.$$
		Độ lệch chuẩn của mẫu số liệu ghép nhóm là $S_1=\sqrt{S_1^2}=\sqrt{7{,}5216}$.\\ 
		Xét mẫu số liệu của cổ  phiếu $B$:\\
		Số trung bình của mẫu  số liệu ghép nhóm là
		$$\overline{x}_2=\dfrac{16\cdot 121+4\cdot  123+3\cdot 125+6\cdot 127+21\cdot 129}{50}=125{,}28.$$
		 Phương sai của mẫu số liệu ghép nhóm là
		$$ S_2^2=\dfrac{1}{50}\left(16\cdot 121^2+4\cdot  123^2+3\cdot 125^2+6\cdot 127^2+21\cdot 129^2\right)-(125{,}48)^2=12{,}4096.$$
		Độ  lệch chuẩn của mẫu số liệu ghép nhóm là $S_2=\sqrt{S_2^2}=\sqrt{12{,}4096}$.\\
		Vậy nếu đánh giá độ rủi ro theo phương sai và độ lệch  chuẩn thì cổ phiếu $A$ có độ rủi ro thấp hơn cổ phiếu $B$.
	} 
\end{bt}

% nguồn: Lop_12_Dai_so_Chuong_3
\begin{bt}%[2D3H2-3]%[ma:b068aa0e]
	Mai và Ngọc cùng sử dụng vòng đeo tay thông minh để ghi lại số bước chân hai bạn đi  mỗi ngày trong một tháng. Kết quả được ghi lại ở bảng sau
	\begin{longtable}{|c|c|c|c|c|c|}
		 \hline
		 Giá trị đại diện & $4$ & $6$ & $8$ & $10$ & $12$\\
		\hline 
		Mai &  $6$ & $7$ & $6$ & $6$ & $5$\\
		 \hline
		Ngọc & $2$ & $5$ & $10$  & $8$ & $2$\\
		 \hline
	\end{longtable}
	\begin{enumerate}
		\item Hãy tính số trung bình và độ lệch chuẩn của mẫu  số liệu ghép nhóm trên.
		\item Nếu so sánh theo  độ lệch chuẩn thì bạn nào có số lượng bước chân đi mỗi ngày đều đặn hơn?
	\end{enumerate}
	 \loigiai{
		\begin{enumerate}
			\item  Cỡ mẫu $ n=30$.\\
			Xét mẫu số liệu  của Mai:\\
			Số trung bình $\overline{x}_1=\dfrac{6\cdot 4+7\cdot 6+6\cdot 8+6\cdot 10+5\cdot  12}{30}=7{,}8$.\\
			Phương sai $S_1^2=\dfrac{\left(6\cdot 4^2+7\cdot 6^2+6\cdot 8^2+6\cdot 10^2+5\cdot 12^2\right)}{30}-7{,}8^2=7{,}56$.\\ 
			Độ lệch chuẩn $s_1=\sqrt{7{,}56}\approx 2{,}75$.\\ 
			Xét mẫu số  liệu của Ngọc:\\
			Số trung bình $\overline{x}_2=\dfrac{2\cdot 4+5\cdot 6+13\cdot  8+8\cdot 10+2\cdot 12}{30}=8{,}2$.\\
			Phương sai $S_2^2=\dfrac{\left(2\cdot 4^2+5\cdot 6^2+13\cdot 8^2+8\cdot 10^2+2\cdot  12^2\right)}{30}-8{,}2^2\approx 3{,}83$.\\
			Độ lệch  chuẩn $s_2=\sqrt{3{,}83}\approx 1{,}96$.
			\item Nếu so sánh theo độ  lệch chuẩn thì bạn Ngọc có số lượng bước chân đi mỗi ngày đều đặn hơn.
		\end{enumerate}
	} 
\end{bt}
